The within-Llama model comparison is also negative. On fixed H, $O_x$ and $E_x$
both have equal-window $L=0.00421052$ and occupied MAE 25.1458 GPU seconds:
neither improves over the other. The registered gate requires at least 20\%
$L$ reduction \emph{and} each window's occupied MAE no worse than $O_x+21$;
passing the latter does not rescue zero improvement. The inherited 20\%
occupied-improvement diagnostic also fails. These locks use development
median/max aggregate service lookups with identical lifecycle endpoints, not
a new Llama demonstration of the Qwen endpoint-adaptation gain. M2 passes on
completed-local common support (minimum 206 requests, maximum MAE 0.02964 s),
but H and guard cover only 51.24\% and 58.90\% of all arrivals on an equal-window
basis. This service-subset result does not establish whole-stream or controller
fidelity; the supplement reports coverage and target disagreement separately.

\subsection{Same-environment lifecycle refit}
W3 asks whether a bounded same-environment refit resolves the original V--E
comparison's environment caveat. It selects $V'$ by loss on two preformal
development canaries; a separate preformal drift probe constrains lifecycle
lower bounds. It freezes the model before CPU replay and evaluates all
36 existing physical records without new GPU runs. On fixed H, equal-window
$L$ falls from 0.01352738 to 0.01284857, a 5.02\% reduction below the registered
10\% threshold, while every window satisfies the occupied-MAE safety line
$V'\leq E+21$. Steady-window $L$ worsens despite the aggregate reduction.
The result is a scientific negative, not evidence that refitting has no effect.
Only lifecycle scalars are eligible for fitting; service, setup, and guard
safety remain fixed. The researchers already knew V3's aggregate results,
although the fitting code never consumed formal outcomes. This is a bounded
retrospective replay, not a newly blinded experiment, and it does not erase the
historical V--E caveat.
